% GENERATED FILE. Edit canonical lessons, then run npm run generate:books.
% canonical-source-hash: fnv1a64:19d6bb82cc7980b3
% canonical-lessons: FR-C312-demenagement, FR-C312-emmenager, FR-C312-cremaillere, FR-C312-locataire, FR-C312-bail

\chapter{Une nouvelle maison}
\label{ch:fr-une-nouvelle-maison}
\hlchaptermodality{312}

\begin{chapteropening}
\textbf{By the end of this chapter:} \emph{I can talk about moving: the move, moving in, the house-warming, the tenant and the lease.}

Complete the last lesson of chapter 312: I can talk about moving: the move, moving in, the house-warming, the tenant and the lease.
\end{chapteropening}

\section[le déménagement]{le déménagement — the move, the house move}

\label{lesson:FR-C312-demenagement}



\begin{quote}
Before the new one: say the French for the tiered wedding cake, then the French for the sugared almond.
\end{quote}

\subsection*{You'll want to know: le déménagement}
\textbf{le déménagement} — \textquotedblleft{}the move, the house move\textquotedblright{}.

\begin{quote}
Le déménagement a pris toute la journée.
\end{quote}

\begin{quote}
\emph{The move took the whole day.}
\end{quote}

The verb is \emph{déménager}.

\begin{cousinweb}
\emph{déménagement} is built on \emph{ménage} \textquotedblleft{}household\textquotedblright{}, which goes back to Vulgar Latin \emph{mansiōnāticum}, built on Latin \emph{mānsiō} \textquotedblleft{}a dwelling\textquotedblright{}, from \emph{manēre} \textquotedblleft{}to stay\textquotedblright{} (Lewis \& Short).
\end{cousinweb}

\begin{grammarlens}[title={What you've built}]
One more word for telling, explaining and describing at B1.
\end{grammarlens}

\subsection*{Your turn}
\begin{itemize}
\raggedright
  \item \emph{Say it:} \emph{le déménagement}
  \item \emph{Say it:} \emph{le déménagement}, once more
  \item \emph{Say it:} a sentence of your own with \emph{le déménagement}
  \item \emph{Recall:} say the French for the tiered wedding cake, then the French for the sugared almond, then say \emph{le déménagement} again
\end{itemize}

\begin{tcolorbox}[breakable,colback=teal!4,colframe=teal!35!black,title={Before you move on}]
What does le déménagement mean? (\textquotedblleft{}The move, the house move\textquotedblright{}.) Say it once more.
\end{tcolorbox}
\section[emménager]{emménager — to move in}

\label{lesson:FR-C312-emmenager}



\begin{quote}
Before the new one: say the French for the sugared almond, then the French for the move.
\end{quote}

\subsection*{You'll want to know: emménager}
\textbf{emménager} — \textquotedblleft{}to move in\textquotedblright{}.

\begin{quote}
Nous avons emménagé dans un nouvel appartement.
\end{quote}

\begin{quote}
\emph{We moved into a new flat.}
\end{quote}

The opposite of \emph{déménager}, to move out.

\begin{cousinweb}
\emph{emménager} is built on \emph{en-} and \emph{ménage} \textquotedblleft{}household\textquotedblright{}, which goes back to Latin \emph{mānsiō} \textquotedblleft{}a dwelling\textquotedblright{} (Lewis \& Short).
\end{cousinweb}

\begin{grammarlens}[title={What you've built}]
One more word for telling, explaining and describing at B1.
\end{grammarlens}

\subsection*{Your turn}
\begin{itemize}
\raggedright
  \item \emph{Say it:} \emph{emménager}
  \item \emph{Say it:} \emph{emménager}, once more
  \item \emph{Say it:} a sentence of your own with \emph{emménager}
  \item \emph{Recall:} say the French for the sugared almond, then the French for the move, then say \emph{emménager} again
\end{itemize}

\begin{tcolorbox}[breakable,colback=teal!4,colframe=teal!35!black,title={Before you move on}]
What does emménager mean? (\textquotedblleft{}To move in\textquotedblright{}.) Say it once more.
\end{tcolorbox}
\section[la crémaillère]{la crémaillère — the house-warming; a pot-hook}

\label{lesson:FR-C312-cremaillere}



\begin{quote}
Before the new one: say the French for the move, then the French for to move in.
\end{quote}

\subsection*{You'll want to know: la crémaillère}
\textbf{la crémaillère} — \textquotedblleft{}the house-warming; a pot-hook\textquotedblright{}.

\begin{quote}
Ils ont pendu la crémaillère samedi soir.
\end{quote}

\begin{quote}
\emph{They had their house-warming on Saturday night.}
\end{quote}

\emph{Pendre la crémaillère} — to throw a house-warming party; once, hanging the pot-hook over the fire was the last step of moving in.

\begin{cousinweb}
\emph{crémaillère} is usually traced to Late Latin \emph{cremaculus} \textquotedblleft{}a pot-hook\textquotedblright{}, from Greek \emph{kremast\'{\={e}}r} \textquotedblleft{}a hanger\textquotedblright{}; the details are not settled (EtymDB lists no parent).
\end{cousinweb}

\begin{grammarlens}[title={What you've built}]
One more word for telling, explaining and describing at B1.
\end{grammarlens}

\subsection*{Your turn}
\begin{itemize}
\raggedright
  \item \emph{Say it:} \emph{la crémaillère}
  \item \emph{Say it:} \emph{la crémaillère}, once more
  \item \emph{Say it:} a sentence of your own with \emph{la crémaillère}
  \item \emph{Recall:} say the French for the move, then the French for to move in, then say \emph{la crémaillère} again
\end{itemize}

\begin{tcolorbox}[breakable,colback=teal!4,colframe=teal!35!black,title={Before you move on}]
What does la crémaillère mean? (\textquotedblleft{}The house-warming; a pot-hook\textquotedblright{}.) Say it once more.
\end{tcolorbox}
\section[le locataire]{le locataire — the tenant}

\label{lesson:FR-C312-locataire}



\begin{quote}
Before the new one: say the French for to move in, then the French for the house-warming.
\end{quote}

\subsection*{You'll want to know: le locataire}
\textbf{le locataire} — \textquotedblleft{}the tenant\textquotedblright{}.

\begin{quote}
Le nouveau locataire est arrivé lundi.
\end{quote}

\begin{quote}
\emph{The new tenant arrived on Monday.}
\end{quote}

The owner is \emph{le propriétaire}.

\begin{cousinweb}
\emph{locataire} is a learned word built on Latin \emph{locātus}, from \emph{locāre} \textquotedblleft{}to place, to let out\textquotedblright{} (Lewis \& Short); EtymDB lists Italian \emph{locatario} beside it.
\end{cousinweb}

\begin{grammarlens}[title={What you've built}]
One more word for telling, explaining and describing at B1.
\end{grammarlens}

\subsection*{Your turn}
\begin{itemize}
\raggedright
  \item \emph{Say it:} \emph{le locataire}
  \item \emph{Say it:} \emph{le locataire}, once more
  \item \emph{Say it:} a sentence of your own with \emph{le locataire}
  \item \emph{Recall:} say the French for to move in, then the French for the house-warming, then say \emph{le locataire} again
\end{itemize}

\begin{tcolorbox}[breakable,colback=teal!4,colframe=teal!35!black,title={Before you move on}]
What does le locataire mean? (\textquotedblleft{}The tenant\textquotedblright{}.) Say it once more.
\end{tcolorbox}
\section[le bail]{le bail — the lease}

\label{lesson:FR-C312-bail}



\begin{quote}
Before the new one: say the French for the house-warming, then the French for the tenant.
\end{quote}

\subsection*{You'll want to know: le bail}
\textbf{le bail} — \textquotedblleft{}the lease\textquotedblright{}.

\begin{quote}
Ils ont signé un bail de trois ans.
\end{quote}

\begin{quote}
\emph{They signed a three-year lease.}
\end{quote}

Plural \emph{les baux}. Informal: \emph{ça fait un bail !} — it has been ages!

\begin{cousinweb}
\emph{bail} is built on \emph{bailler} \textquotedblleft{}to hand over\textquotedblright{} (EtymDB), which is inherited from Latin \emph{bāiulāre} \textquotedblleft{}to carry a load\textquotedblright{} (Lewis \& Short).
\end{cousinweb}

\begin{grammarlens}[title={What you've built}]
One more word for telling, explaining and describing at B1.
\end{grammarlens}

\subsection*{Your turn}
\begin{itemize}
\raggedright
  \item \emph{Say it:} \emph{le bail}
  \item \emph{Say it:} \emph{le bail}, once more
  \item \emph{Say it:} a sentence of your own with \emph{le bail}
  \item \emph{Recall:} say the French for the house-warming, then the French for the tenant, then say \emph{le bail} again
\end{itemize}

\begin{tcolorbox}[breakable,colback=teal!4,colframe=teal!35!black,title={Before you move on}]
What does le bail mean? (\textquotedblleft{}The lease\textquotedblright{}.) Say it once more.
\end{tcolorbox}
